\begin{table}[H]
    \footnotesize
    \centering
    \begin{threeparttable}
        \caption{Journal LLM tonal composites, comparisons to \textit{AER}}
        \label{table3llm_journal}
        \begin{tabular}{p{3cm}S@{}S@{}S@{}S@{}S@{}}
            \toprule
            &{\crcell[b]{G1: Creativity\\[-0.1cm]\& Hedging}}&{\crcell[b]{G2: Assert.\\[-0.1cm]\& Voice}}&{\crcell[b]{G3: Structural\\[-0.1cm]Directness}}&{\crcell[b]{G4: Stance\\[-0.1cm]\& Novelty}}&{\crcell[b]{G5: Support\\[-0.1cm]\& Impact}}\\
            \midrule
            \textit{Econometrica}         &       -0.73***&       -1.23***&       -1.87***&       -0.40*  &        0.07   \\
                                          &      (0.21)   &      (0.32)   &      (0.23)   &      (0.21)   &      (0.25)   \\
            \textit{JPE}                  &       -2.50***&       -0.99***&       -2.72***&        0.33   &        0.79***\\
                                          &      (0.21)   &      (0.31)   &      (0.23)   &      (0.21)   &      (0.25)   \\
            \textit{QJE}                  &       -0.48***&        0.00   &       -0.08   &        0.13** &        0.39***\\
                                          &      (0.10)   &      (0.13)   &      (0.09)   &      (0.06)   &      (0.11)   \\
            \midrule
            No. observations              &       9,117   &       9,117   &       9,117   &       9,117   &       9,117   \\
            \bottomrule
        \end{tabular}
        \begin{tablenotes}
            \tiny
            \item \textit{Notes}. Figures are the estimated coefficients on the journal dummy variables from column (2) in~\autoref{table3llm_FemRatio}. Each column corresponds to a separate OLS regression of LLM tonal composite scores on the ratio of female authors, controlling for editor effects.
        \end{tablenotes}
    \end{threeparttable}
\end{table}